
\documentclass[11pt,a4paper]{article}

% --- Packages ---
\usepackage[utf8]{inputenc}
\usepackage[T1]{fontenc}
\usepackage{amsmath}
\usepackage{amsfonts}
\usepackage{amssymb}
\usepackage{graphicx}
\usepackage{hyperref}
\usepackage{geometry}
\usepackage{booktabs}
\usepackage{listings}
\usepackage{xcolor}

% --- Margins ---
\geometry{margin=1in}

% --- Formatting ---
\hypersetup{
    colorlinks=true,
    linkcolor=blue,
    filecolor=magenta,      
    urlcolor=cyan,
    pdftitle={SAAS User Manual},
}

\definecolor{codegray}{rgb}{0.5,0.5,0.5}
\definecolor{codeblue}{rgb}{0.1,0.1,0.8}
\lstset{
    backgroundcolor=\color{gray!10},
    basicstyle=\ttfamily\small,
    breaklines=true,
    captionpos=b,
    commentstyle=\color{codegray},
    keywordstyle=\color{codeblue},
}

% --- Title ---
\title{Software for the Analysis of Atomic Spectra (SAAS) \\ \Large User Manual}
\author{Gillian Nave, Jacob Ward, Christian Clear}
\date{\today}

\begin{document}

\maketitle

\tableofcontents
\newpage

\section{Introduction}
The Software for the Analysis of Atomic Spectra (SAAS) is a unified suite designed for experimental atomic physics. It replaces a diverse set of legacy scripts (notably \texttt{pybranch.py}) with a modern, modular application written in Python 3.

SAAS provides a structured environment for managing spectroscopic data, identifying spectral lines, and calculating fundmental atomic properties such as **Branching Fractions (BFs)** and **Transition Probabilities (Einstein A-coefficients)**.

\section{Core Architecture}
\subsection{HDF5 Project Management}
Unlike legacy workflows that rely on scattered ASCII files, SAAS utilizes a single **HDF5 (.h5)** project file. This format acts as a container for:
\begin{itemize}
    \item Raw and calibrated spectra (NSO and FITS formats).
    \item Experimental linelists (.lin and .txt).
    \item Master energy level tables and theoretical calculations.
    \item Versioned analysis results and metadata.
\end{itemize}

\subsection{The Main Interface}
The main window consists of an **HDF5 Tree Browser** on the left and a **Data Preview/Metadata Pane** on the right.
\begin{description}
    \item[Tree Selection:] Click any item to see a quick preview (first 200 rows) or a spectral plot.
    \item[Metadata View:] The second tab displays all HDF5 attributes associated with the selected object.
    \item[Full Table Viewer:] Right-click any dataset in the tree and select \textbf{Full Table}. This opens a searchable, filterable window containing the entire dataset, essential for navigating large experimental linelists.
\end{description}

\section{Installation and Build}
\subsection{System Requirements}
SAAS is designed for Linux environments (GNOME/Wayland). It requires:
\begin{itemize}
    \item Python 3.8 or higher.
    \item Dependencies: \texttt{pandas, numpy, h5py, matplotlib, PyQt5, click, astropy}.
\end{itemize}

\subsection{Building the Executable}
A build script \texttt{build.sh} is provided to create a standalone binary:
\begin{lstlisting}[language=bash]
chmod +x build.sh
./build.sh
\end{lstlisting}
This packages the code and libraries into a single file in the \texttt{dist/} directory and generates a Linux \texttt{.desktop} file for menu integration.

\section{Data Import}
SAAS provides dedicated wizards to ensure data integrity and schema compliance.

\subsection{Spectra and Linelists}
\begin{itemize}
    \item \textbf{Main Spectrum:} Imports .raw/.dat files with their .hdr headers.
    \item \textbf{Calibration Spectrum:} Associates a standard lamp spectrum with an existing main spectrum to facilitate intensity calibration.
    \item \textbf{Binary Linelists:} Directly parses Xgremlin \texttt{.lin} files, automatically applying damping parameter transformations.
\end{itemize}

\subsection{Generic Table Wizard}
For Energy Levels and Previous Identifications, the **Import Wizard** allows users to:
\begin{enumerate}
    \item Preview ASCII/CSV files.
    \item Assign columns to specific fields (e.g., mapping a column to \texttt{lifetime\_unc\_frac}).
    \item Store source metadata (author, date, notes) directly in the project file.
\end{enumerate}

\section{Interactive Branching Fraction Analysis}
The analysis module (accessible via \texttt{Ctrl+R}) is the heart of SAAS.

\subsection{The Analysis Workflow}
\begin{enumerate}
    \item \textbf{Level Selection:} Choose a master levels file and select an upper level of interest. SAAS filters the Previous Identifications list to show transitions originating from that level.
    \item \textbf{Data Selection:} Select the experimental linelists for comparison. \textit{Note: Checking a linelist automatically enables its corresponding raw spectrum for plotting.}
    \item \textbf{Normalization:} Select a strong line observed in most spectra. Right-click and select \textbf{Normalize to Reference} to set its intensity to 1000 and rescale all other lines.
    \item \textbf{Transfer Calibration:} If a spectrum lacks the primary reference line, right-click a common ``Transfer Line'' to bridge the intensity scales between spectra.
\end{enumerate}

\subsection{Outlier Highlighting}
SAAS performs real-time outlier detection. If an individual intensity measurement deviates by more than \textbf{3-sigma} from the weighted mean, the cell is highlighted in **light red**. Sigma is calculated as the root-sum-of-squares of the absolute individual uncertainty and the mean uncertainty.

\section{Physical Calculations}
\subsection{Uncertainty Model}
SAAS implements the model established by Sikström et al. (2002). The total absolute uncertainty of an individual measurement is derived from:
\begin{enumerate}
    \item \textbf{Statistical Uncertainty ($U_{stat}$):} 
    \begin{equation}
        U_{stat} = \frac{1.5}{SNR \cdot \sqrt{n}}
    \end{equation}
    where $n = \text{FWHM} / \text{resolution}$.
    \item \textbf{Calibration Uncertainty ($Unc_{cal}$):} This term accounts for the increasing error as one moves away from the spectrum's strongest calibrated line ($\sigma_{maxI}$):
    \begin{equation}
        Unc_{cal} = \text{CalUncPer1000} \cdot \frac{|\sigma - \sigma_{maxI}|}{1000}
    \end{equation}
\end{enumerate}

\subsection{Weighted Mean and Constraints}
Lines are weighted by $1/Unc_{total}^2$. To prevent strong lines from disproportionately biasing the results, uncertainties are clipped to a minimum of 6\% (a weight cap of 555).

\subsection{Residual Estimation}
Unobserved decay channels are accounted for by a ``residual'' fraction calculated from theoretical transition probabilities ($A_{theo}$):
\begin{equation}
    \text{Fraction}_{\text{resid}} = \frac{\sum A_{\text{unobserved}} \cdot \tau_{\text{lifetime}}}{1000}
\end{equation}
A default uncertainty of 50\% is assigned to the residual.

\section{Reporting and Export}
\subsection{The Rule of 20}
To maintain scientific rigor, results are rounded according to the metrological **Rule of 20**:
\begin{itemize}
    \item If the first significant digit of the absolute uncertainty is \textbf{1}, the value is rounded to two significant digits of the uncertainty.
    \item If the first significant digit is \textbf{2 or greater}, the value is rounded to one significant digit of the uncertainty.
\end{itemize}

\subsection{Export Formats}
\begin{description}
    \item[LaTeX:] Generates a \texttt{tabular} environment with all percentage signs and special characters properly escaped for direct insertion into papers.
    \item[MRT (Machine Readable Table):] Exports using the \texttt{astropy.io.ascii.mrt} library. This format is the standard for AAS journals (ApJ, AJ). SAAS provides both **Truncated** (matching the UI) and **Raw** (full precision) MRT export options.
\end{description}

\section{Troubleshooting}
\begin{itemize}
    \item \textbf{Visual Artifacts:} If random colored patterns appear on Linux, SAAS automatically forces the \texttt{xcb} (X11) backend and software rendering (\texttt{LIBGL\_ALWAYS\_SOFTWARE=1}) to ensure stability on Wayland.
    \item \textbf{File Access:} If HDF5 files on network drives show locking errors, run: \\ \texttt{export HDF5\_USE\_FILE\_LOCKING=FALSE} in your shell.
\end{itemize}

\end{document}
